% !TITLE 燕歌行·其二
% !AUTHOR 曹丕
% !DYNASTY 魏晋
% !GENRE 七言古诗
% !ORDER 6
% !EDITION 《先秦汉魏晋南北朝诗》，逯钦立辑校，中华书局
% !EDITION_STATUS pending

\begin{poem}
别日何易会日难，山川悠远路漫漫。
郁陶\textsuperscript{1}思君未敢言，寄声\textsuperscript{2}浮云往不还。
涕零雨面\textsuperscript{3}毁容颜，谁能怀忧独不叹？
展诗\textsuperscript{4}清歌聊自宽，乐往哀来摧肺肝。
耿耿\textsuperscript{5}伏枕不能眠，披衣出户步东西，仰看星月观云间。
飞鸧\textsuperscript{6}晨鸣声可怜\textsuperscript{7}，留连顾怀不能存\textsuperscript{8}。
\end{poem}

\begin{annotations}
\notetext{1}{郁陶：郁陶（yù yáo），忧思郁积、闷在心里的样子。}
\notetext{2}{寄声：托人捎话。这里指托浮云带去口信。}
\notetext{3}{涕零雨面：泪水像雨一样落满脸上。}
\notetext{4}{展诗：展开诗卷。}
\notetext{5}{耿耿：心中有事、不能安宁的样子。}
\notetext{6}{鸧：鸧（cāng），鸟名。}
\notetext{7}{可怜：可哀，惹人怜惜。}
\notetext{8}{留连顾怀不能存：徘徊眷恋，心神不能安定。存，安、自持。}
\end{annotations}

\begin{appreciation}
《燕歌行》共两首，这是第二首。前一首从秋风写起，一路写到牵牛织女隔河相望；这一首没有秋景的铺垫，一开口便是离别。

“别日何易会日难，山川悠远路漫漫。”分别的日子来得何其容易，重逢的日子却那样难；山川悠远，长路漫漫。离别只需一天，重逢却要隔着整片山河，一易一难之间，说尽了人间最寻常的不公平。六百多年后李商隐写“相见时难别亦难”，又翻进一层，连离别也说成了难事。

“郁陶思君未敢言，寄声浮云往不还。”思念郁积在心里，她却不敢说出口，只好托天上的浮云捎去口信，可浮云飘走了，再也没有回来。前一首说“不敢忘”，这一首说“未敢言”，忘不掉，又说不出，曹丕笔下的女子总是把最深的情绪压在一个“敢”字底下。与其说这是闺中女子的羞怯，不如说是写诗人的性情。

“涕零雨面毁容颜，谁能怀忧独不叹？”泪水如雨，流了满面，容颜都憔悴了；心怀忧愁，谁又能独自忍住不叹息呢？这一问像是在替自己开脱：我也想忍住，可是忍不住。克制的人偶尔失守，反倒更叫人心疼。

“展诗清歌聊自宽，乐往哀来摧肺肝。”她展开诗卷，清声歌唱，想姑且宽慰自己；可歌声里那一点快乐刚刚过去，悲哀便又涌上来，摧折着肺肝。前一首弹琴，“短歌微吟不能长”；这一首唱歌，“乐往哀来”。两首诗里的女子都想借音乐排遣忧愁，两次都没有成功。曹丕大约很明白，诗与歌能让人暂时忘记痛苦，却替人消除不了痛苦。

“耿耿伏枕不能眠，披衣出户步东西，仰看星月观云间。”心事耿耿，伏在枕上睡不着，便披上衣服走出门去，来回踱步，抬头看星月，看云。曹丕在《杂诗》里也写过一个失眠的人：“展转不能寐，披衣起彷徨。彷徨忽已久，白露沾我裳。”同样是披衣，同样是彷徨，诗里的女子和写诗的帝王，在许多个长夜里做着同一件事。

“飞鸧晨鸣声可怜，留连顾怀不能存。”天亮了，一只鸧鸟在晨光里鸣叫，叫声凄楚可怜；她徘徊着、眷恋着，心神再也无法安定。前一首停在“夜未央”，长夜还远没有尽头；这一首一直写到了清晨。从夜到晓，她一夜未眠，所有的辗转反侧，只换来一声鸟鸣。

两首合读，像同一个女子的两个夜晚，又像同一个夜晚的前后两半。一首仰望天上的星河，一首走到了户外；一首还要问牛郎织女何辜，一首连问也不问了，只是听着鸟叫。曹丕没有给她一个结局，远人没有回来，书信也没有来，天亮了，日子还要过下去。

鸟还在叫。她站在晨光里，什么也没有说。
\end{appreciation}
